```latex
\documentclass[11pt,a4paper]{article}

% ---------- Packages ----------
\usepackage[utf8]{inputenc}
\usepackage[T1]{fontenc}
\usepackage[english]{babel}
\usepackage{amsmath, amssymb}
\usepackage{graphicx}
\usepackage{booktabs}
\usepackage{siunitx}
\usepackage{geometry}
\usepackage{hyperref}
\usepackage{caption}
\usepackage{subcaption}

\geometry{margin=2.5cm}

\hypersetup{
    colorlinks=true,
    linkcolor=blue,
    citecolor=blue,
    urlcolor=blue
}

% ---------- Title ----------
\title{\textbf{Coarsening and Arrest in Colloidal Gels:}\\
Modelling via Brownian Dynamics}

\author{Bérénice Derai \\ 66034}

\date{5 June 2026}

\begin{document}

\maketitle

\begin{abstract}
This project explores the dynamics of colloidal gels, focusing on the processes of coarsening and arrest. A minimal model of attractive Brownian particles is constructed and their dynamics are simulated using an Euler--Maruyama scheme with a Lennard-Jones pair potential. The growth of the structure factor $S(q,t)$ is analysed by comparing two regimes: a strongly attractive system, $\varepsilon = 5 k_B T$, exhibiting kinetic arrest, and a weakly attractive liquid reference, $\varepsilon = 1.5 k_B T$, that undergoes continuous coarsening.

An ageing study is additionally performed: the mean-squared displacement (MSD) is computed for multiple waiting times $t_w$, revealing $t_w$-dependent dynamics in the gel, associated with non-ergodic ageing, and $t_w$-independent diffusion in the liquid, associated with ergodic behaviour. Several numerical challenges, including integrator instability, inadequate simulation time, and an unphysical reference state, were identified and corrected iteratively. The theoretical framework draws on the works of Doi for statistical mechanics and colloidal diffusion, and on the works of van Saarloos \textit{et al.} on ageing and non-equilibrium states.
\end{abstract}

\tableofcontents
\newpage

\section{Introduction}

Colloidal suspensions are systems composed of particles with sizes typically between $1\,\mathrm{nm}$ and $1\,\mu\mathrm{m}$ dispersed in a solvent. They provide a useful model system for the physics of non-equilibrium condensed matter because their dynamics can be directly related to thermal fluctuations, interparticle interactions, and collective structural evolution.

When sufficiently strong attractive interactions are induced between colloidal particles, for example by the addition of salt in a charged suspension, by polymer depletion, or by unscreened van der Waals interactions, the system may undergo phase separation towards a dense state. However, the intrinsically slow dynamics of massive colloidal particles can trap the system in a non-equilibrium state before phase separation is complete. This phenomenon is known as gelation by arrested phase separation.

The aim of this work is to construct a minimal model of attractive Brownian particles, to simulate its dynamics using an Euler--Maruyama scheme, and to analyse the growth of the structure factor $S(q,t)$ in the presence and absence of gelation. An ageing study based on the mean-squared displacement is also carried out in order to probe the non-ergodic character of the arrested gel.

The report is organised as follows. First, the theoretical background of Brownian motion, colloidal diffusion, and attractive interactions is introduced. The mechanism of phase separation and gelation by kinetic arrest is then discussed. The structure factor is presented as the central observable used to characterise coarsening and dynamic scaling. Finally, the numerical implementation and simulation results are analysed, with a particular focus on the comparison between a gel regime and a weakly attractive liquid reference.
```
```latex
\section{Brownian Motion and Colloidal Diffusion}

\subsection{The Langevin Equation}

Following Doi, the motion of a colloidal particle suspended in a solvent can be described by separating the force exerted by the fluid into two contributions: a deterministic viscous friction force and a stochastic random force arising from microscopic collisions with solvent molecules.

The equation of motion is the Langevin equation

\begin{equation}
m\ddot{\mathbf r}
=
-\zeta \dot{\mathbf r}
-\nabla U
+
\mathbf F_r(t),
\label{eq:langevin}
\end{equation}

where $-\zeta \dot{\mathbf r}$ is the viscous friction force, $U(\mathbf r)$ is the interaction potential, and $\mathbf F_r(t)$ is the stochastic random force.

For a spherical colloid of radius $a$ immersed in a solvent of viscosity $\eta$, the Stokes friction coefficient is

\begin{equation}
\zeta = 6\pi \eta a.
\end{equation}

The random force satisfies

\begin{equation}
\langle \mathbf F_r(t)\rangle = 0,
\end{equation}

and obeys the fluctuation--dissipation theorem

\begin{equation}
\left\langle
F_{r,\mu}(t)
F_{r,\nu}(t')
\right\rangle
=
2\zeta k_B T
\delta_{\mu\nu}
\delta(t-t').
\end{equation}

This relation ensures that thermal fluctuations compensate viscous dissipation and maintain thermal equilibrium at temperature $T$.

\subsection{Overdamped Limit: Brownian Dynamics}

For colloidal particles, the inertial relaxation time

\begin{equation}
\tau_m = \frac{m}{\zeta}
\end{equation}

is typically of order $10^{-8}\,\mathrm{s}$, whereas the diffusive timescale

\begin{equation}
\tau_D = \frac{a^2}{D}
\end{equation}

is of order $10^{-3}\,\mathrm{s}$.

Since $\tau_m \ll \tau_D$, inertia relaxes much faster than diffusion and the inertial term can be neglected. The Langevin equation therefore reduces to

\begin{equation}
\dot{\mathbf r}
=
-\frac{D}{k_B T}\nabla U
+
\frac{1}{\zeta}\mathbf F_r(t),
\label{eq:bd}
\end{equation}

where the Einstein relation

\begin{equation}
D=\frac{k_B T}{\zeta}
\end{equation}

has been used.

Equation~(\ref{eq:bd}) clearly separates deterministic drift induced by interactions from stochastic diffusion generated by thermal agitation. This overdamped description forms the basis of Brownian Dynamics simulations.

\section{Colloidal Interactions}

\subsection{Pair Potentials Between Colloids}

Two particles of radius $a$ cannot overlap. This excluded-volume constraint is represented by a hard-core interaction, which is often replaced in simulations by a steep but continuous repulsive potential.

In addition to steric repulsion, colloidal particles interact through attractive forces. Quantum fluctuations of electronic dipoles generate long-range van der Waals attractions between dielectric materials. For charged particles dispersed in an electrolyte, the DLVO theory describes the competition between electrostatic repulsion and van der Waals attraction.

Increasing the salt concentration screens electrostatic interactions, lowers the repulsive barrier, and can render the effective interaction attractive at short distances. This mechanism is one of the most common routes to colloidal gelation.

\subsection{Choice of Potential for Simulations}

The simulations performed in this work use a truncated and shifted Lennard--Jones potential,

\begin{equation}
U_{\mathrm{LJ}}(r)
=
4\varepsilon
\left[
\left(\frac{\sigma}{r}\right)^{12}
-
\left(\frac{\sigma}{r}\right)^6
\right]
-
U_{\mathrm{LJ}}(r_{\mathrm{cut}}),
\label{eq:lj}
\end{equation}

with a cutoff distance

\begin{equation}
r_{\mathrm{cut}} = 2.5\sigma.
\end{equation}

The Lennard--Jones potential was chosen because of its widespread use as a benchmark model for attractive colloidal systems and its numerical simplicity.

Two regimes are considered:

\begin{enumerate}
\item \textbf{Gel regime}

\[
\varepsilon = 5\,k_B T,
\qquad
r_{\mathrm{cut}} = 2.5\sigma.
\]

The deep attractive well promotes cluster formation and kinetic arrest.

\item \textbf{Liquid reference regime}

\[
\varepsilon = 1.5\,k_B T,
\qquad
r_{\mathrm{cut}} = 2.5\sigma.
\]

The attraction is sufficiently weak that the system remains ergodic and continues to coarsen throughout the simulation.
\end{enumerate}
```
```latex
\section{Phase Separation in Colloidal Suspensions}

\subsection{Phase Diagram and Spinodal Instability}

The phase diagram of a system of attractive particles exhibits a liquid--gas-like coexistence region. For sufficiently strong attractions, the homogeneous state becomes unstable with respect to density fluctuations and the system separates into particle-rich and particle-poor regions.

The stability of the homogeneous state is determined by the curvature of the free-energy density $f(\phi)$, where $\phi$ denotes the volume fraction. When

\begin{equation}
\frac{\partial^2 f}{\partial \phi^2} < 0,
\end{equation}

the homogeneous phase is unstable. This region of the phase diagram is known as the spinodal region.

Inside the spinodal region, arbitrarily small density fluctuations grow spontaneously without the need for nucleation, leading to spinodal decomposition.

\subsection{Spinodal Decomposition: the Cahn--Hilliard Theory}

A simple description of spinodal decomposition is provided by the Cahn--Hilliard theory. The free-energy functional is written as

\begin{equation}
F[\phi]
=
\int
\left[
f(\phi)
+
\frac{\kappa}{2}
|\nabla\phi|^2
\right]
dV,
\end{equation}

where the gradient term penalises the formation of interfaces.

Considering small fluctuations around a homogeneous state

\begin{equation}
\phi(\mathbf r,t)
=
\phi_0+\delta\phi(\mathbf r,t),
\end{equation}

and linearising the dynamics yields the growth rate of a Fourier mode of wavevector $q$,

\begin{equation}
\sigma(q)
=
M_0 q^2
\left[
|f''(\phi_0)|
-
\kappa q^2
\right].
\label{eq:cahn}
\end{equation}

When $f''(\phi_0)<0$, long-wavelength modes become unstable and grow exponentially.

Only modes satisfying

\begin{equation}
q<q_c
=
\sqrt{\frac{|f''(\phi_0)|}{\kappa}}
\end{equation}

are unstable.

The most rapidly growing mode corresponds to

\begin{equation}
q^*
=
\frac{q_c}{\sqrt{2}},
\end{equation}

which defines an initial characteristic length scale

\begin{equation}
\ell^*
=
\frac{2\pi}{q^*}.
\end{equation}

This selected length scale is the origin of the peak observed in the structure factor during phase separation.

\section{Gelation as Arrested Phase Separation}

\subsection{General Mechanism}

A widely accepted description of colloidal gelation interprets the gel state as an arrested phase separation.

Following a quench into the spinodal region, density fluctuations grow and particle-rich domains begin to form through spinodal decomposition. As phase separation progresses, the local volume fraction inside the dense regions increases continuously.

When the local density becomes sufficiently high, particle motion slows dramatically. Rearrangements become increasingly difficult and the dense phase eventually undergoes a glass transition.

At this point the coarsening process is interrupted before complete phase separation is achieved. Domain growth stops and the system becomes trapped in a non-equilibrium state composed of a percolating network of dynamically arrested particles.

The resulting arrested structure is referred to as a colloidal gel.

\subsection{Ageing and Non-Equilibrium States}

Unlike equilibrium systems, colloidal gels continue to evolve slowly after their preparation.

Their properties depend explicitly on the waiting time $t_w$ elapsed since the quench. Such systems are said to exhibit ageing.

A commonly used observable is the mean-squared displacement (MSD),

\begin{equation}
\mathrm{MSD}(t,t_w)
=
\left\langle
\left|
\mathbf r(t_w+t)-\mathbf r(t_w)
\right|^2
\right\rangle.
\end{equation}

In a fully arrested gel, particles become temporarily trapped by their neighbours and the MSD develops an intermediate plateau,

\begin{equation}
\mathrm{MSD}(t)
\rightarrow
r_{\mathrm{cage}}^2,
\end{equation}

where $r_{\mathrm{cage}}$ is the typical cage size.

Furthermore, the MSD becomes dependent on the waiting time $t_w$, reflecting the progressive slowing down of the dynamics as the system ages.

By contrast, an ergodic liquid exhibits waiting-time-independent dynamics and the MSD is independent of $t_w$.

\section{Structure Factor, Coarsening and Dynamic Scaling}

\subsection{Time-Dependent Structure Factor}

The central observable of this work is the time-dependent structure factor.

For a system of $N$ particles located at positions $\mathbf r_j(t)$, the microscopic density can be written as

\begin{equation}
\rho(\mathbf r,t)
=
\sum_{j=1}^{N}
\delta(\mathbf r-\mathbf r_j(t)).
\end{equation}

Its Fourier transform is

\begin{equation}
\rho_{\mathbf q}(t)
=
\sum_{j=1}^{N}
e^{i\mathbf q\cdot\mathbf r_j(t)}.
\end{equation}

The isotropically averaged structure factor is then defined as

\begin{equation}
S(q,t)
=
\frac{1}{N}
\left|
\sum_{j=1}^{N}
e^{i\mathbf q\cdot\mathbf r_j(t)}
\right|^2.
\label{eq:Sq}
\end{equation}

The structure factor measures density fluctuations at different length scales.

As phase separation develops, a peak emerges at a characteristic wavevector $q^*(t)$.

The corresponding characteristic domain size is

\begin{equation}
\ell(t)
=
\frac{2\pi}{q^*(t)}.
\label{eq:length}
\end{equation}

Monitoring the evolution of $q^*(t)$ therefore provides direct information on the growth of domains during coarsening.

\subsection{Diffusive Coarsening and Dynamic Scaling}

For diffusion-controlled phase separation, the characteristic length scale follows the Lifshitz--Slyozov law

\begin{equation}
\ell(t)
\propto
t^{1/3}.
\end{equation}

As coarsening proceeds, the morphology becomes statistically self-similar: the system appears identical when lengths are rescaled by $\ell(t)$.

This dynamic scaling hypothesis implies that the structure factor obeys

\begin{equation}
S(q,t)
=
\ell(t)^d
\,\widetilde S
\bigl(q\ell(t)\bigr),
\label{eq:scaling}
\end{equation}

where $d$ is the spatial dimension and $\widetilde S$ is a universal scaling function.

When dynamic scaling holds, structure-factor curves measured at different times collapse onto a single master curve when plotted as

\[
S(q,t)/\ell(t)^d
\qquad \text{versus} \qquad
q\ell(t).
\]

In a gelled system, kinetic arrest interrupts coarsening and this scaling behaviour eventually breaks down, signalling the emergence of a frozen characteristic length scale.
``` 
\end{document}
